\begin{abstract}
Deep learning models deployed in safety-critical settings can fail silently under distribution shift, producing confident but unreliable outputs with no way to self-diagnose the failure. We study this in seismic phase picking, the task of identifying earthquake wave arrivals (`picks') in seismograms. PhaseNet, a widely-deployed deep learning model trained on Northern California data, sees its P-wave detection rate collapse from 83.9\% to 35.1\% when deployed to Kazakhstan, a seismically active but data-sparse region where automated tools are most needed yet least validated. PhaseNet's own confidence scores for each pick are nearly indistinguishable between the two regions (AUROC of 0.699 in California vs 0.701 in Kazakhstan), so a monitoring agency relying on confidence alone would have no way of knowing that detection performance had collapsed. We develop a regime-aware deferral framework to address this. The framework combines an external Hidden Markov Model of the seismic environment with Beta-Bernoulli estimation of how often PhaseNet's picks are correct in each regime. The regime-aware score defers 90.9\% of picks made during seismically active periods, compared to 36.4\% when deferring on PhaseNet's confidence alone, directing scarce analyst time toward regimes where reliability rests on the fewest observations, at a marginal cost in accuracy, the fraction of picks handled correctly (0.619 vs 0.651). The regions that stand to benefit most from deploying this framework are precisely those where automated tools have not been tested against local ground truth, making deferral mechanisms that know when not to trust the model essential for ensuring that data-sparse regions receive the same monitoring safeguards as data-rich ones.
\end{abstract}
